% Source: https://artofproblemsolving.com/community/c3h1578767_rational_functions
% Stable identity: aops-topic-1578767-q01
Suppose $f(x)=mx+c$ for all $x\ne2$. Compute three values:
\[
f(0)=\dfrac{-3}{-2}=\dfrac{3}{2},\qquad f(1)=\dfrac{-2}{-1}=2,\qquad f(3)=\dfrac{0}{1}=0 .
\]
From $f(0)$ and $f(1)$ we get $c=\dfrac{3}{2}$ and $m=2-\dfrac{3}{2}=\dfrac{1}{2}$, so we would need $f(3)=\dfrac{3}{2}+\dfrac{3}{2}=3$. But $f(3)=0$, a contradiction. So $f$ is not linear. $\blacksquare$

Another way to see it: $f(x)=1-\dfrac{1}{x-2}$, whose graph is a hyperbola with the vertical asymptote $x=2$; a linear function is defined for every $x$ and has no asymptote. The "rate of change" between two points is not constant: between $x=0$ and $x=1$ it is $\dfrac{1}{2}$, between $x=1$ and $x=3$ it is $-1$.
